\chapter{武汉大学硕士学位论文印制规定}

(原文网址: \url{http://www.gs.whu.edu.cn/index.php/index-view-aid-7930.html}.)

为了规范硕士学位论文印制，使之符合国家标准局编写格式，及时向社会提供查阅，促进国内外学术交流，我校对硕士学位论文印制格式作出如下规定：

\textbf{一、论文装订的排列顺序}

（一）封面（学术型硕士格式见附件一，专业硕士格式见附件二）

（二）论文英文题目（专用—页纸，上方为题目用宋体2号字，下方为研究生姓名宋体4号字、外文专业应有中文题目）

（三）论文原创性声明（格式和要求见附件三）

（四）中文摘要

（五）英文摘要

（六）目录

（七）引言（绪论）

（八）正文

（九）中外文参考文献（格式和要求见附件四）

（十）附录

（十一）致谢或后记

\textbf{二、论文印制规格及要求}

（一）论文封面统一用120克铜版纸，封面底色为白色

（二）论文用A4纸（210×297mm）标准大小的白纸，\colorbox{yellow}{必须双面印制}，前置部分例外。

（三）论文在印制时，纸张四周留足空白边缘，即：每页上方（天头）、下方（地脚）、左侧（订口）、右侧（切口）应分别留出25mm以上的空白边缘。

\textbf{三、论文封面格式}

1、论文封面（封一）的内容包括：

（一）分类号。必须在封面左上角注明分类号，一般应注明《中国图书资料分类法》的类号，同时尽可能注明《国际十进分类法UDC》的类号。

（二）密级。如涉密论文，密级及保密年限由学校相关保密部门审核认定后按国家规定的保密条例在右上角注明密级（如系公开型论文则可不注明密级）。

（三）编号。武汉大学编号为10486，标注在封面右上角。

（四）论文题目。题目必须用楷体标准一号字标注于明显的位置，应是集中概括论文最重要的内容，一般不超过20个字，以有助于选定关键词和编制题录。题目不能用缩略词，首字母缩写字、字符、代号和公式等，题目语意未尽，可用副标题补充说明。外语专业的论文题目一般采用英文，英文题目不宜超过10个实词。

（五）论文作者姓名。

（六）论文作者学号。

（七）指导教师姓名、职称。指导教师姓名必须是填写被我校批准招收博士生、硕士生的教师。

（八）专业名称。学术型硕士专业名称必须是我校已有学位授予权的学科专业，并按国家颁布的学科专业目录中二级学科、专业名称印制。专业学位硕士填写专业学位类别，下设领域的具体到具体领域。

（九）论文完成时间。用汉字书写，只写年月，如“二〇一六年四月”。

2、论文封面的格式，请严格按“标准样本”（学术型硕士详见附件一，专业型硕士详见附件二）制作。

\textbf{四、论文英文题目}

论文英文题目专用一页纸，“英文题目”用Times new Roman字体二号字，其下“研究生姓名”用Times new Roman字体四号字；外语专业应为中文题目。

\textbf{五、论文原创性声明}（见附件三）

论文原创性声明用黑体小二号字，内容用宋体四号字。

\textbf{六、中文摘要}

中文摘要用黑体小二号字，内容用宋体小四号字，页码用罗马数字单独编排，并标注在每页页脚中部。

\textbf{七、英文摘要}

英文摘要用加粗Times New Roman小二号字，内容用Times New Roman小四号字，页码续接中文摘要的页码。

\textbf{八、论文关键词}

每篇论文必须选取3-5个中、英文关键词，排在其论文摘要的左下方，用黑体小四号字。

\textbf{九、目录}

目录是论文的提纲，也是论文组成部分的小标题。排列顺序是：1.中文摘要；2.英文摘要；3.引言（绪论）；4.正文章节；5.参考文献；6.致谢或后记。

\textbf{十、引言（绪论）}

论文的页码由引言（绪论）的首页开始，作为第1页，并为右页，一律用阿拉伯数字连续编排页码，必须统一标注在每页页脚中部。

\textbf{十一、正文}

正文是学位论文的核心部分，必须由另页开始，一级标题之间换页，二级标题之间空行；内容一律用宋体小四号字，字间距设置为标准字间距，行间距设置为最小值20磅，各章、节应有序号。

\textbf{十二、参考文献}（“文后参考文献著录格式”见附件四）

参考文献用黑体四号字，内容用宋体五号字。

\textbf{十三、附录}

附录是作为论文主体部分的补充项目，并不是必需的。



\textbf{十四、论文制作时须加页眉}


页眉从中文摘要开始至论文末，偶数页码内容为：武汉大学硕士学位论文，奇数页码内容为学位论文题目。